% ch1-gls.tex
\chapter{A Complete Worked Example: Generalized Least Squares}
\label{ch:gls}

\begin{answer}
ulm.lam turns the formula \(\beta=(X\T M\inv X)\inv X\T M\inv y\) into
eleven calls: three copies, two Cholesky factorizations, a triangular solve
with many right-hand sides, a symmetric rank-\(k\) update, a matrix-vector
product and three triangular solves with one right-hand side. No inverse is
formed, \(M\) is factorized once although it is inverted twice, and the
matrix \(L\inv X\) is computed once and used twice. The facts that make this
possible (\(M\) is SPD, \(X\) has full rank, \(X\T M\inv X\) is SPD) are
carried in MLIR tensor types; two of them are declared and one is inferred.
The two factorizations are the only checks of these facts in the generated
code.
\end{answer}

This chapter follows one program through every representation.
\Cref{fig:pipeline} shows the stages; each section below explains one of
them before the next one appears. All listings are the actual output of the
commands in \cref{app:repro} (files in \code{outputs/gls/}).

\begin{figure}[t]
\centering
\begin{tikzpicture}[
  font=\sffamily\small,
  node distance=0.95cm,
  rep/.style={draw=accent,fill=accent!6,rounded corners=2pt,align=left,
    text width=5.6cm,inner sep=4pt},
  fact/.style={align=left,text width=4.6cm,font=\sffamily\scriptsize,text=metacolor,
    anchor=west},
  pass/.style={font=\sffamily\scriptsize,text=accent!70!black,align=left,anchor=west,
    inner sep=1pt},
  arr/.style={-{Stealth[length=2mm]},draw=accent!80!black,thick}]
  \node[rep] (jl) {\textbf{Julia}: \code{@lam function gls(...)}};
  \node[rep,below=of jl] (tr) {\textbf{ulm\_lam} on tensors: \code{mul}, \code{transpose},
    \code{inverse}\\\scriptsize\code{1-traced.mlir}};
  \node[rep,below=of tr] (no) {\textbf{ulm\_lam}, normal form\\\scriptsize\code{2-normal.mlir}};
  \node[rep,below=of no] (in) {\textbf{ulm\_lam}, properties inferred\\\scriptsize\code{3-inferred.mlir}};
  \node[rep,below=of in] (se) {\textbf{ulm\_lamk} on tensors (DPS):
    \code{potrf}, \code{trsm}, \code{syrk}, \ldots\\\scriptsize\code{5-selected.mlir}};
  \node[rep,below=of se] (bu) {\textbf{ulm\_lamk} on memrefs, workspace, \code{info}\\\scriptsize\code{6-bufferized.mlir}};
  \node[rep,below=of bu] (cc) {\textbf{EmitC} \(\to\) \textbf{C++} calling BLAS/LAPACK\\\scriptsize\code{7-gls.cpp}, compiled to a shared library};
  \foreach \a/\b/\t in {jl/tr/{tracing (operator overloading)},
                       tr/no/{\code{ulm-lam-normalize}},
                       no/in/{\code{ulm-lam-infer-properties}},
                       in/se/{\code{loop-invariant-code-motion}, \code{ulm-lam-select-kernels}},
                       se/bu/{\code{strip-properties}, \code{one-shot-bufferize}, \code{ulm-lam-workspace}, \code{ulm-lam-info}},
                       bu/cc/{\code{convert-ulm-lamk-to-ulmblas}, \code{mlir-translate}}}
    \draw[arr] (\a.south) -- node[pass,xshift=3pt] {\t} (\b.north);
  \node[fact] at ([xshift=0.35cm]jl.east) {facts in Julia types: \code{Mat(:n, :m, :full_rank)}};
  \node[fact] at ([xshift=0.35cm]tr.east) {facts in tensor types: shape, \code{#ulm_lam.props}; sizes in \code{ulm_lam.dims}};
  \node[fact] at ([xshift=0.35cm]in.east) {stronger types on results: \code{spd} for \(X\T M\inv X\)};
  \node[fact] at ([xshift=0.35cm]se.east) {facts turned into kernel choices and flags (\code{uplo}, \code{trans})};
  \node[fact] at ([xshift=0.35cm]bu.east) {encodings gone; buffers, copies, offsets; \code{info} codes};
  \node[fact] at ([xshift=0.35cm]cc.east) {pointers, strides, \code{size_t} sizes; no types beyond \code{double *}};
\end{tikzpicture}
\caption{The stages of the GLS example. Left: the representation and the
pass that produces it. Right: where the facts about the matrices live at
that stage.}
\label{fig:pipeline}
\end{figure}

% ---------------------------------------------------------------------------
\section{The problem and its assumptions}
\label{sec:gls-problem}

In the linear model \(y = X\beta + e\) with an error \(e\) of covariance
\(\sigma^2 M\), the GLS estimate of \(\beta\in\R^m\) is
\[
  \beta = (X\T M\inv X)\inv X\T M\inv y,
  \qquad X\in\R^{n\times m},\; M\in\R^{n\times n},\; y\in\R^n .
\]
The formula needs two facts:
\begin{enumerate}
\item \(M\) is SPD, so \(M\inv\) exists and is SPD;
\item \(X\) has full column rank (\(\operatorname{rank} X = m \le n\)).
\end{enumerate}
Together they make \(S = X\T M\inv X\) SPD: for \(v \ne 0\), \(Xv \ne 0\)
by (2), and \(v\T S v = (Xv)\T M\inv (Xv) > 0\) by (1). So \(S\) is
invertible and \(\beta\) is well defined.

A numerical analyst does not form either inverse. With the Cholesky
factorization \(M = L_1 L_1\T\) (\(L_1\) lower triangular):
\[
  W = L_1\inv X,\quad z = L_1\inv y,\quad S = W\T W = L_2 L_2\T,\quad
  \beta = L_2\invT L_2\inv\, W\T z .
\]
\(W\) appears twice; \(S\) is symmetric, so only one triangle is computed;
\(S\) is SPD, so it is factorized by Cholesky as well. This is what ulm.lam
derives, and the rest of the chapter shows how.

% ---------------------------------------------------------------------------
\section{The Julia frontend code}
\label{sec:gls-julia}

The program, as in the README and in \code{examples/gls.jl}:

\begin{lstlisting}[style=jl]
using UlmLam, LinearAlgebra

@lam function gls(X::Mat(:n, :m, :full_rank), M::Mat(:n, :n, :spd), y::Vec(:n))
    inv(X' * inv(M) * X) * X' * inv(M) * y
end

beta = gls(X, M, y)    # Julia arrays; the first call compiles
\end{lstlisting}

The body is the formula. The parameter types are not Julia types but
ulm.lam types, values of the Julia structs \code{MatType}
(\srcref{julia/UlmLam/src/frontend.jl}{39}{45}):

\begin{itemize}
\item \code{Mat(:n, :m, :full_rank)} is a matrix with \(n\) rows and \(m\)
  columns and the property \code{full_rank}. The symbols \code{:n} and
  \code{:m} are \emph{run-time sizes}: their values are taken from the
  arrays at each call, and equal symbols mean equal sizes. A number
  instead of a symbol is a size fixed at compile time.
\item \code{Mat(:n, :n, :spd)} is an \(n\times n\) SPD matrix.
\item \code{Vec(:n)} is a vector, represented as an \(n\times 1\) matrix.
\end{itemize}

The macro \code{@lam} does not compile anything. It wraps the function so
that the first call with Julia arrays traces the body, compiles the traced
program into a shared library (about half a second), and calls it.

% ---------------------------------------------------------------------------
\section{Tracing: from Julia to the first IR}
\label{sec:gls-trace}

\begin{steps}
\item[Input] the Julia function, called once with placeholder values of type \code{LamValue}.
\item[Representation] each Julia operation on a \code{LamValue} appends one line of MLIR text.
\item[Rule] operator overloading: \code{*} emits \code{ulm_lam.mul}, \code{'} emits \code{ulm_lam.transpose}, \code{inv} emits \code{ulm_lam.inverse}.
\item[Output] a \code{func.func} in the \code{ulm_lam} dialect (\texttt{1-traced.mlir}).
\item[Theory] a tracing front end of an embedded DSL (domain-specific language), \FG{3.7}{Embedded DSLs and tracing/AST-capturing front ends}.
\end{steps}

The overloaded product (\srcref{julia/UlmLam/src/frontend.jl}{222}{231})
already checks the inner sizes, by comparing sizes that may be symbols:

\begin{lstlisting}[style=jl]
function Base.:*(a::LamValue, b::LamValue)
    ...
    ta, tb = mat(a, "*"), mat(b, "*")
    ta.cols == tb.rows ||
        throw(LamError("product of $ta and $tb: $(ta.cols) columns but " *
                       "$(tb.rows) rows"))
    op!("ulm_lam.mul", [a, b], newtype(ta.rows, tb.cols))
end
\end{lstlisting}

\paragraph{The signature} The traced function begins (one line in the
file, wrapped here):

\begin{lstlisting}
func.func @gls(%s.n: index, %s.m: index,
    %p.X: tensor<?x?xf64, #ulm_lam.props<full_rank>> {ulm_lam.in, ulm_lam.dims = [1, 0]},
    %p.M: tensor<?x?xf64, #ulm_lam.props<spd>> {ulm_lam.in, ulm_lam.dims = [0, 0]},
    %p.y: tensor<?x1xf64> {ulm_lam.in, ulm_lam.dims = [0]})
    -> (tensor<?x1xf64> {ulm_lam.out, ulm_lam.dims = [1]}) {
\end{lstlisting}

Read it from left to right:
\begin{itemize}
\item \code{\%s.n} and \code{\%s.m} are the run-time sizes, as
  \code{index} parameters of the function. The generated C function receives
  them as \code{size_t} arguments.
\item \code{\%p.X} has type \code{tensor<?x?xf64, #ulm_lam.props<full_rank>>}:
  a rank-2 tensor of \code{f64} with two dynamic dimensions (\code{?}) and,
  in the \emph{encoding} slot of the builtin tensor type, the property set
  \code{full_rank}.
\item \code{ulm_lam.dims = [1, 0]} is an argument attribute. It says which
  \code{index} parameter each dynamic dimension is: the rows of \code{\%p.X}
  are parameter 1 (\(m\)), its columns parameter 0 (\(n\)). So \code{\%p.X} is
  \(m\times n\), not \(n\times m\): Julia stores matrices column-major, the
  ulm.lam ABI (application binary interface) is row-major, and the memory of
  the Julia matrix \(X\) \emph{is} the row-major matrix \(X\T\). The frontend
  types the parameter as \(X\T\) (\code{abitype},
  \srcref{julia/UlmLam/src/frontend.jl}{100}{106}) and hands the traced
  function \code{transpose(\%p.X)}, which is \(X\). No data is copied for this.
\item \code{ulm_lam.in} marks a read-only argument, \code{ulm_lam.out} the
  result, which the caller allocates.
\end{itemize}

\paragraph{The body} Twelve operations, one per Julia operation:

\begin{lstlisting}
%v1 = ulm_lam.transpose %p.X  : ... -> tensor<?x?xf64>        // X   (Julia's X)
%v2 = ulm_lam.transpose %p.M  : ... -> tensor<?x?xf64>        // M
%v3 = ulm_lam.transpose %v1   : ...                           // X'
%v4 = ulm_lam.inverse %v2     : ...                           // inv(M)
%v5 = ulm_lam.mul %v3, %v4    : ...                           // X' * inv(M)
%v6 = ulm_lam.mul %v5, %v1    : ...                           // ... * X
%v7 = ulm_lam.inverse %v6     : ...                           // inv(...)
%v8 = ulm_lam.transpose %v1   : ...                           // X'
%v9 = ulm_lam.inverse %v2     : ...                           // inv(M), again
%v10 = ulm_lam.mul %v7, %v8   : ...
%v11 = ulm_lam.mul %v10, %v9  : ...
%v12 = ulm_lam.mul %v11, %p.y : ... -> tensor<?x1xf64>
return %v12 : tensor<?x1xf64>
\end{lstlisting}

(The comments are added here; the types after \code{:} are shortened.)
Every intermediate result has the plain type \code{tensor<?x?xf64>}: the
frontend writes no properties for results. It also records exactly what
the Julia code does: Julia's left-to-right \code{*} gives binary products
nested to the left, and \code{inv(M)} appears twice (\code{\%v4},
\code{\%v9}) because the formula contains it twice.

\begin{theory}[ (\FGs{3.7})]
ulm.lam is an embedded DSL with a \emph{tracing} front end: the host
language runs the program on symbolic stand-ins and records the operations,
as JAX does with Python. Host-language control flow is evaluated during
tracing and disappears; only the recorded operations reach the compiler.
What is known when (Julia values at trace time, sizes at call time, entries
at run time) is a question of binding times, the sixth thread of the Field
Guide's conceptual map, and of \FG{8.6}{Multi-stage programming and
binding-time analysis}.
\end{theory}

% ---------------------------------------------------------------------------
\section{What the types encode, and what they do not}
\label{sec:gls-types}

Before the first pass runs, the IR carries everything the compiler will
know about the problem. \Cref{tab:gls-info} lists where each piece of
information is.

\begin{table}[t]
\centering\small
\begin{tabularx}{\linewidth}{@{}P{3.3cm}P{4.4cm}L@{}}
\toprule
\textbf{Information} & \textbf{Where in the IR} & \textbf{Origin} \\
\midrule
element type \code{f64} & tensor element type & \Odecl{} by the call
  (Float64 arrays; Float32 arrays compile an \code{f32} version) \\
rank 2 & tensor rank; ODS constraint \code{2DTensorOf<[F32, F64]>} & \Over{}
  by the generated operand checks \\
static sizes & tensor shape & \Over{} by op verifiers when static \\
which run-time size & \code{index} parameters and the argument attribute
  \code{ulm_lam.dims} & \Odecl{} by the frontend; read by an \Oana{}
  (\code{dimOf}, \cref{sec:shapes-dims}) \\
matrix properties & tensor encoding \code{#ulm_lam.props<...>} & \Odecl{} on
  parameters; \Oinf{} on results (\cref{sec:gls-infer}) \\
may the function overwrite it? & argument attributes \code{ulm_lam.in} /
  \code{ulm_lam.inout} & \Odecl{}; drives bufferization
  (\cref{sec:gls-storage}) \\
the entries & nowhere (run-time values) & checked only indirectly, by the
  factorizations \Orun{} \\
\bottomrule
\end{tabularx}
\caption{Where the GLS program's information lives in the first IR.}
\label{tab:gls-info}
\end{table}

\paragraph{Why the properties are a tensor encoding} The builtin
\code{RankedTensorType} has an optional \emph{encoding} attribute. ulm.lam
defines an attribute, \code{#ulm_lam.props<...>}, and places it there,
instead of defining its own matrix type. Upstream passes and One-Shot
Bufferize then see ordinary tensors, and the properties are removed in one
step (\code{ulm-lam-strip-properties}) once they are no longer needed. The
attribute is a bit set over fifteen named facts; its declaration is in
\srcref{include/ulm_lam/UlmLamAttrs.td}{1}{75} and is discussed in
\cref{sec:infra-props}.

\paragraph{What the type of \code{\%p.M} says} \code{tensor<?x?xf64,
#ulm_lam.props<spd>>} reads ``some \(n\times n\) matrix that is SPD''. The
compiler takes this as given. Nothing in the IR proves that the caller's
array is symmetric or positive definite, and the compiler never reads the
entries. This is the first boundary of the system: \Odecl{} facts enter at
the parameters (and at \code{ulm_lam.assume}), and everything else is
derived from them.

\begin{theory}[ (\FGs{1.7})]
A type that restricts a base type by a predicate,
\(\{A : \R^{n\times n} \mid A \text{ is SPD}\}\), is a \emph{refinement
type}, \FG{1.7}{Refinement types}. ulm.lam's predicates come from a fixed
set of fifteen atoms and their conjunctions, like the qualifiers of liquid
types~\cite{liquid,refinement}. \Cref{ch:infra,ch:facts} make the comparison precise, including
what is missing (no proof obligations at call sites, no predicates that
relate a matrix to other values).
\end{theory}

% ---------------------------------------------------------------------------
\section{Normal form}
\label{sec:gls-normal}

\begin{steps}
\item[Input] the traced IR of \cref{sec:gls-trace}.
\item[Representation] the same ops; patterns are C++ classes derived from \code{OpRewritePattern}.
\item[Rule] fold \(A\T{}\T \to A\); fold \(A\T \to A\) when \(A\) is symmetric; flatten nested products into one n-ary \code{ulm_lam.mul} (if the inner product has no other use).
\item[Output] \texttt{2-normal.mlir}, below.
\item[Theory] greedy pattern rewriting to a fixpoint, \FG{3.4}{Pattern rewriting: greedy drivers, canonicalization, PDL}; termination, \FG{3.5}{Term rewriting theory: confluence, termination, strategies}.
\end{steps}

The pass \code{ulm-lam-normalize} collects the canonicalization patterns of
the ops and runs MLIR's greedy driver (\code{applyPatternsGreedily},
\srcref{lib/UlmLamPasses.cpp}{105}{133}). The output (types shortened):

\begin{lstlisting}
%0 = ulm_lam.transpose %arg2 : <X^T, full_rank> -> tensor<?x?xf64>             // X
%1 = ulm_lam.inverse %arg3 : <M, spd> -> tensor<?x?xf64>                       // M^-1
%2 = ulm_lam.mul %arg2, %1, %0 : ... -> tensor<?x?xf64>                       // X^T M^-1 X
%3 = ulm_lam.inverse %2 : tensor<?x?xf64> -> tensor<?x?xf64>
%4 = ulm_lam.inverse %arg3 : <M, spd> -> tensor<?x?xf64>                       // M^-1
%5 = ulm_lam.mul %3, %arg2, %4, %arg4 : ... -> tensor<?x1xf64>                // S^-1 X^T M^-1 y
return %5 : tensor<?x1xf64>
\end{lstlisting}

Three rules did the work:
\begin{itemize}
\item \code{FoldTranspose} (\srcref{lib/UlmLamDialect.cpp}{604}{623})
  replaced \(M\T\) by \(M\). It may do so only because the \emph{type} of
  \code{\%p.M} says \code{spd}, which implies \code{symmetric}. A rewrite
  that depends on a property in a type: this is the first place where a
  declared fact changes the program.
\item The same pattern folded \((X\T)\T\) to the parameter itself, so
  \(X\T\) is now \code{\%arg2} and \(X\) is \code{\%0}.
\item \code{Flatten} (\srcref{lib/UlmLamDialect.cpp}{578}{602}) merged the
  left-nested binary products into one three-factor and one four-factor
  \code{ulm_lam.mul}. An n-ary product leaves the parenthesization open;
  choosing it is the job of kernel selection.
\end{itemize}

Note what did \emph{not} happen: \code{\%1} and \code{\%4} are still two
\code{inverse} ops of the same value. The normalization pipeline contains
no CSE (common subexpression elimination). The factorization will be shared
anyway, by a memo table in kernel selection (\cref{sec:gls-select}).

\begin{theory}[ (\FGs{3.4}, \FGs{3.5})]
Each pattern is a rewrite rule \(\ell \Rightarrow r\) with a side
condition (``\(A\) is symmetric'', ``the inner product has one use''). The
greedy driver applies rules until none matches, a fixpoint of the rewrite
relation. That the fixpoint is reached is a termination argument (each rule
decreases a well-founded measure, as the pass description in
\srcref{include/ulm_lam/UlmLamPasses.td}{8}{19} states); whether it is
unique is confluence. Both are questions of
\FG{3.5}{Term rewriting theory: confluence, termination, strategies}.
\end{theory}

% ---------------------------------------------------------------------------
\section{Inferring properties}
\label{sec:gls-infer}

\begin{steps}
\item[Input] the normal form.
\item[Representation] result types without properties (\code{tensor<?x?xf64>}).
\item[Rule] for every op in program order, compute the result type from the operand types (\code{InferTypeOpInterface::inferReturnTypes}) and replace the old one.
\item[Output] \texttt{3-inferred.mlir}: \(M\inv\) is SPD, \(X\T M\inv X\) is SPD.
\item[Theory] type inference in synthesis mode, \FG{1.10}{Type inference beyond Hindley–Milner: bidirectional and constraint-based}; the property set as an abstract domain, \FG{5.2}{Abstract interpretation: Galois connections, widening, narrowing}.
\end{steps}

The pass \code{ulm-lam-infer-properties}
(\srcref{lib/UlmLamPasses.cpp}{139}{183}) changes only types. The
interesting lines of its output:

\begin{lstlisting}
%1 = ulm_lam.inverse %arg3 : tensor<?x?xf64, #ulm_lam.props<spd>>
       -> tensor<?x?xf64, #ulm_lam.props<square, symmetric, spd, non_singular, full_rank>>
%2 = ulm_lam.mul %arg2, %1, %0 : ... -> tensor<?x?xf64, #ulm_lam.props<spd>>
%3 = ulm_lam.inverse %2 : tensor<?x?xf64, #ulm_lam.props<spd>>
       -> tensor<?x?xf64, #ulm_lam.props<square, symmetric, spd, non_singular, full_rank>>
\end{lstlisting}

\paragraph{\(M\inv\) is SPD} The rule for \code{inverse} keeps the
properties that survive inversion (\code{inverseProperties}, a mask). The
list printed for \code{\%1} is the \emph{closure} of \code{spd}: SPD
implies symmetric, square and non-singular, and square plus non-singular
implies full rank (\code{closure}, \srcref{lib/UlmLamDialect.cpp}{97}{123}).

\paragraph{\(X\T M\inv X\) is SPD} The rule for \code{mul} recognizes
products that are their own transpose. Its core
(\srcref{lib/UlmLamDialect.cpp}{238}{269}):

\begin{lstlisting}[style=cpp]
/// Products that read the same backwards when transposed,
/// F_1 ... F_k with F_i^T = F_(k+1-i):
///   Z^T Z      (k even)            spsd, spd if Z has full column rank;
///   Z^T M Z    (k odd, M middle)   symmetric if M is, spsd if M is, spd if
///                                  M is spd and Z has full column rank.
P palindromeProperties(ValueRange f, ArrayRef<P> ps)
{
    size_t k = f.size(), h = k / 2;
    ...
    for (size_t i = 0; i < h; ++i) {
        auto [a, ta] = leafOf(f[i]);           // (matrix, transposed?)
        auto [b, tb] = leafOf(f[k - 1 - i]);
        if (a != b || (ta == tb && !isSymmetric(a)))
            return P::general;
    }
    ...
    P m = ps[h];                               // the middle factor
    if (has(m, P::spd) && fullColumnRank(z, pz))
        return P::spd;
    ...
\end{lstlisting}

For \code{mul \%arg2, \%1, \%0}: the outer factors are \code{\%arg2}
(\(X\T\), the parameter) and \code{\%0} (\(X\), a transpose of
\code{\%arg2}); \code{leafOf} reduces both to \code{\%arg2} with opposite
transposition, so the product has the form \(Z\T M Z\). The middle factor
\code{\%1} is SPD. \code{fullColumnRank} accepts \(Z = X\) because \(X\) is
\code{full_rank} \emph{and its sizes are dynamic}: the source comment
(\srcref{lib/UlmLamDialect.cpp}{222}{236}) says that a full-rank matrix of
unknown shape is taken to have full column rank, since the formula inverts
\(X\T M\inv X\) and needs it. Hence \code{spd}.

\Needspace{9\baselineskip}
\paragraph{Where each fact came from}
\begin{center}\small
\begin{tabularx}{\linewidth}{@{}lL@{}}
\toprule
\(M\) is SPD & \Odecl{} (\code{Mat(:n, :n, :spd)}) \\
\(X\) has full rank & \Odecl{} (\code{Mat(:n, :m, :full_rank)}) \\
\(n \ge m\), i.e.\ full \emph{column} rank & \Onone{} at compile time: taken for granted because the sizes are dynamic; \Orun{} indirectly, if \(S\) is singular its Cholesky factorization fails \\
\(M\inv\) is SPD & \Oinf{} from \(M\) SPD \\
\(X\T M\inv X\) is SPD & \Oinf{} from the three facts above, by the palindrome rule \\
\bottomrule
\end{tabularx}
\end{center}

The inferred fact is exactly as reliable as the facts it was inferred from.
It is a theorem about real matrices, \emph{conditional on} the declarations.

\begin{theory}[ (\FGs{1.10}, \FGs{5.1}, \FGs{5.2})]
Computing result types bottom-up from operand types is type
\emph{synthesis} in the sense of bidirectional typing,
\FG{1.10}{Type inference beyond Hindley–Milner: bidirectional and
constraint-based}. There are no type variables and no unification: one walk
in program order suffices because SSA values are defined before they are
used and the IR here has no cycles. The property sets with
``absence means unknown'' form a lattice ordered by implication, and each
op's rule is a sound transfer function on it: \FG{5.1}{Lattices, monotone
frameworks and fixpoints (a sharper recap)} and \FG{5.2}{Abstract
interpretation: Galois connections, widening, narrowing}.
\Cref{sec:facts-lattice} develops both views.
\end{theory}

% ---------------------------------------------------------------------------
\section{Kernel selection}
\label{sec:gls-select}

\begin{steps}
\item[Input] the inferred IR: two products with inverses, with properties on every value.
\item[Representation] each product becomes a sequence of \emph{leaves}: matrices (possibly transposed) and triangular solves.
\item[Rule] factorize each inverted matrix once, by its properties; parenthesize each product by a dynamic program over kernels and flop counts; reuse a sub-product computed earlier, also transposed.
\item[Output] \code{ulm_lamk} kernels on tensors (\texttt{5-selected.mlir}).
\item[Theory] instruction selection by tiling with dynamic programming, \FG{4.5}{Instruction selection: tree/DAG tiling, SelectionDAG vs GlobalISel}; progressive lowering, \FG{3.3}{Progressive lowering and the dialect conversion framework}.
\end{steps}

The pass \code{ulm-lam-select-kernels} (class \code{Selector},
\srcref{lib/UlmLamPasses.cpp}{318}{2265}) does three things for GLS.

\paragraph{1. Factorizations from properties} Every inverted matrix is
factorized once, directly after its definition (\code{factorise},
\srcref{lib/UlmLamPasses.cpp}{855}{903}). The choice is read off the type:
triangular: no factorization; \code{spd}: Cholesky (\code{potrf}, LAPACK's
Cholesky factorization of a positive definite matrix); otherwise: LU
factorization with partial pivoting (\code{getrf}). Results are memoized per
matrix, so \code{\%1} and \code{\%4}, both \code{inverse \%arg3}, share one
\code{potrf}. An inverse becomes solves: \(M\inv = L_1\invT L_1\inv\).

\paragraph{2. The first product} \(X\T M\inv X\) becomes the leaves
\(X\T,\ L_1\invT,\ L_1\inv,\ X\). A dynamic program over intervals of
leaves (\code{solveChain}, \srcref{lib/UlmLamPasses.cpp}{1059}{1138})
considers, for each interval, a \code{gemm} of two sub-intervals, a solve
applied from the left or right in place, and the \emph{Gram} form \(Z\T Z\)
when the left half is the transpose of the right half. Here the halves are
\(X\T L_1\invT\) and \(L_1\inv X\), transposes of each other, so the
product is \(W\T W\) with \(W = L_1\inv X\): one \code{trsm} (triangular
solve with many right-hand sides, BLAS level 3, \(n^2 m\) flops) and one
\code{syrk} (symmetric rank-\(k\) update, one triangle, \(m^2 n\) flops).
The result \(S\) is SPD by its type, so it is factorized by \code{potrf}.

\paragraph{3. The second product} \(S\inv X\T M\inv y\) becomes
\(L_2\invT,\ L_2\inv,\ X\T,\ L_1\invT,\ L_1\inv,\ y\). The interval
\(X\T L_1\invT\) is the transpose of \(W\), which the first product
computed; the table of computed products is keyed by leaves and
transposition, so it costs nothing and \(W\) is reused. What remains is
\(z = L_1\inv y\) (\code{trsv}, triangular solve with one right-hand side),
\(W\T z\) (\code{gemv}, matrix-vector product), and two \code{trsv} with
\(L_2\).

\paragraph{Output} After \code{ulm-lam-strip-properties} and
\code{canonicalize} (\texttt{5-selected.mlir}; constants omitted):

\begin{lstlisting}
%result, %info = ulm_lamk.potrf uplo(lower) outs(%arg3 : tensor<?x?xf64>) -> tensor<?x?xf64>
%0 = tensor.empty(%arg1, %arg1) : tensor<?x?xf64>
%1 = tensor.empty(%arg0, %arg1) : tensor<?x?xf64>
%2 = ulm_lamk.gecopy trans(trans) ins(%arg2 : tensor<?x?xf64>) outs(%1 : tensor<?x?xf64>) -> tensor<?x?xf64>
%3 = ulm_lamk.trsm side(left) uplo(lower) trans(none) diag(non_unit) alpha(%cst_0)
       ins(%result : tensor<?x?xf64>) outs(%2 : tensor<?x?xf64>) -> tensor<?x?xf64>
%4 = ulm_lamk.syrk uplo(lower) trans(trans) alpha(%cst_0) beta(%cst)
       ins(%3 : tensor<?x?xf64>) outs(%0 : tensor<?x?xf64>) -> tensor<?x?xf64>
%result_1, %info_2 = ulm_lamk.potrf uplo(lower) outs(%4 : tensor<?x?xf64>) -> tensor<?x?xf64>
%5 = ulm_lamk.trsv uplo(lower) trans(none) diag(non_unit)
       ins(%result : tensor<?x?xf64>) outs(%arg4 : tensor<?x1xf64>) -> tensor<?x1xf64>
%6 = tensor.empty(%arg1) : tensor<?x1xf64>
%7 = ulm_lamk.gemv trans(trans) alpha(%cst_0) beta(%cst)
       ins(%3, %5 : tensor<?x?xf64>, tensor<?x1xf64>) outs(%6 : tensor<?x1xf64>) -> tensor<?x1xf64>
%8 = ulm_lamk.trsv uplo(lower) trans(none) diag(non_unit)
       ins(%result_1 : tensor<?x?xf64>) outs(%7 : tensor<?x1xf64>) -> tensor<?x1xf64>
%9 = ulm_lamk.trsv uplo(lower) trans(trans) diag(non_unit)
       ins(%result_1 : tensor<?x?xf64>) outs(%8 : tensor<?x1xf64>) -> tensor<?x1xf64>
return %9 : tensor<?x1xf64>
\end{lstlisting}

Each kernel is in destination-passing style (DPS)~\cite{dps}: the operand that BLAS
overwrites is the \code{outs} operand, and on tensors the op returns the
new value of it. \code{potrf} ``overwrites'' \code{\%arg3}, which is the
caller's \(M\); on tensors that is harmless (a new value is returned), and
bufferization will decide what it means for memory. \code{\%3}, the matrix
\(W\), has two uses (\code{syrk} and \code{gemv}): one value, used twice.
\code{gecopy trans(trans)} makes \(X\) (the transpose of \code{\%arg2})
into a fresh row-major matrix that \code{trsm} can overwrite.

Julia's \code{@code_lam} prints the same sequence in the user's terms:

\begin{textout}
gls(X::n×m full_rank, M::n×n spd, y::n)
  potrf    L1 ← chol(M)                      M = L1 L1ᵀ, Cholesky
  copy     T1 ← X
  trsm     T1 ← L1⁻¹ T1
  syrk     T2 ← T1ᵀ T1                       one triangle only
  potrf    L2 ← chol(T2)                     T2 = L2 L2ᵀ, Cholesky
  trsv     T3 ← L1⁻¹ y
  gemv     T4 ← T1ᵀ T3
  trsv     T4 ← L2⁻¹ T4
  trsv     T4 ← L2⁻ᵀ T4
  return T4
if a factorisation fails, the call throws InfoError with
  info 1:  Cholesky factorisation of M: not positive definite
  info 2:  Cholesky factorisation of a computed matrix: not positive definite
\end{textout}

\begin{figure}[t]
\centering
\begin{tikzpicture}[font=\sffamily\small,
  v/.style={draw=accent,fill=accent!6,rounded corners=2pt,minimum height=0.75cm,
    align=center,inner sep=3pt},
  inp/.style={v,fill=white},
  k/.style={font=\sffamily\scriptsize,text=accent!70!black},
  arr/.style={-{Stealth[length=1.8mm]},draw=accent!80!black,semithick}]
  \node[inp] (X) at (0,0) {\(X\)};
  \node[inp] (M) at (3.6,0) {\(M\)};
  \node[inp] (y) at (7.2,0) {\(y\)};
  \node[v,fill=oinf!12,draw=oinf] (W) at (0,-1.5) {\(W = L_1\inv X\)};
  \node[v] (L1) at (3.6,-1.5) {\(L_1\)};
  \node[v] (z) at (7.2,-1.5) {\(z = L_1\inv y\)};
  \node[v] (S) at (0,-3) {\(S = W\T W\)};
  \node[v] (r) at (7.2,-3) {\(r = W\T z\)};
  \node[v] (L2) at (0,-4.5) {\(L_2\)};
  \node[v] (b) at (7.2,-4.5) {\(\beta = L_2\invT L_2\inv r\)};
  \draw[arr] (X) -- node[k,right] {copy, trsm} (W);
  \draw[arr] (M) -- node[k,right] {potrf} (L1);
  \draw[arr] (y) -- node[k,right] {copy, trsv} (z);
  \draw[arr] (L1) -- (W);
  \draw[arr] (L1) -- (z);
  \draw[arr] (W) -- node[k,right] {syrk} (S);
  \draw[arr] (S) -- node[k,right] {potrf} (L2);
  \draw[arr,draw=oinf,very thick] (W.south east) -- (r.north west);
  \draw[arr] (z) -- node[k,right] {gemv} (r);
  \draw[arr] (r) -- node[k,right] {trsv, trsv} (b);
  \draw[arr] (L2) -- (b);
\end{tikzpicture}
\caption{Values and dependences of the selected kernels. \(W\) (green) has
two uses; \(L_1\) has two uses. Each value is computed once.}
\label{fig:gls-dag}
\end{figure}

\begin{theory}[ (\FGs{4.5}, \FGs{3.3})]
Kernel selection is instruction selection with BLAS and LAPACK as the
instruction set: leaves are operands, kernels are tiles that cover an
interval of the product, and the dynamic program picks the cheapest tiling,
as in tree tiling with dynamic programming,
\FG{4.5}{Instruction selection: tree/DAG tiling, SelectionDAG vs
GlobalISel}. The interval structure is that of the matrix chain problem
(generalized to solves and transposes by Barthels et al.~\cite{gmc}). The
pass replaces one dialect by another in a single step and keeps the
structure needed later (DPS, flags), which is \emph{progressive lowering},
\FG{3.3}{Progressive lowering and the dialect conversion framework}.
It does \emph{not} use MLIR's dialect conversion framework: it is a custom
C++ traversal (\cref{sec:kernels-select}).
\end{theory}

\begin{limit}
The selected sequence equals the formula \emph{over the real numbers}. In
floating point it is a different computation with different rounding:
\(W\T W\) is not bitwise \((X\T M\inv) X\), and its accuracy depends on the
condition number of \(W\) squared. \Cref{sec:gls-numbers} shows a
difference in the last bit; \cref{sec:kernels-fp} shows a case where the
two differ completely.
\end{limit}

% ---------------------------------------------------------------------------
\section{From values to memory: bufferization, workspace and \texttt{info}}
\label{sec:gls-storage}

\begin{steps}
\item[Input] kernels on tensors in DPS; read-only arguments marked \code{ulm_lam.in}.
\item[Representation] after bufferization: kernels on \code{memref}s, with no results.
\item[Rule] One-Shot Bufferize writes a kernel's result into the buffer of its \code{outs} operand unless that buffer is read-only or still read later; then it copies first.
\item[Output] \texttt{6-bufferized.mlir}: three copies, one workspace, an \code{info} result.
\item[Theory] \FG{10.3}{Bufferization and destination-passing style}; \FG{2.7}{Value semantics vs reference/memory semantics in IRs}.
\end{steps}

The core of the bufferized function (the workspace protocol at its start
and the \code{scf.if} nesting of the \code{info} checks are left out; see
\cref{sec:storage-workspace} for both):

\begin{lstlisting}
%view = memref.view %25[%c0][%arg0, %arg0] : memref<?xi8> to memref<?x?xf64>
memref.copy %arg3, %view : memref<?x?xf64> to memref<?x?xf64>          // copy of M
%info = ulm_lamk.potrf uplo(lower) outs(%view : memref<?x?xf64>)        // L1 in the copy
...                                                                     // if info == 0:
%view_1 = memref.view %25[%4][%arg1, %arg1] : memref<?xi8> to memref<?x?xf64>
%view_2 = memref.view %25[%9][%arg0, %arg1] : memref<?xi8> to memref<?x?xf64>
ulm_lamk.gecopy trans(trans) ins(%arg2 : memref<?x?xf64>) outs(%view_2 : memref<?x?xf64>)
ulm_lamk.trsm side(left) uplo(lower) trans(none) diag(non_unit) alpha(%cst)
    ins(%view : memref<?x?xf64>) outs(%view_2 : memref<?x?xf64>)       // W
ulm_lamk.syrk uplo(lower) trans(trans) alpha(%cst) beta(%cst_0)
    ins(%view_2 : memref<?x?xf64>) outs(%view_1 : memref<?x?xf64>)     // S
%info_3 = ulm_lamk.potrf uplo(lower) outs(%view_1 : memref<?x?xf64>)   // L2
...                                                                     // if info_3 == 0:
%view_4 = memref.view %25[%13][%arg0] : memref<?xi8> to memref<?x1xf64>
memref.copy %arg4, %view_4 : memref<?x1xf64> to memref<?x1xf64>        // copy of y
ulm_lamk.trsv ... ins(%view : memref<?x?xf64>) outs(%view_4 : memref<?x1xf64>)        // z
ulm_lamk.gemv trans(trans) ... ins(%view_2, %view_4 : ...) outs(%arg5 : memref<?x1xf64>)
ulm_lamk.trsv uplo(lower) trans(none) ... ins(%view_1 : ...) outs(%arg5 : memref<?x1xf64>)
ulm_lamk.trsv uplo(lower) trans(trans) ... ins(%view_1 : ...) outs(%arg5 : memref<?x1xf64>)
\end{lstlisting}

\paragraph{Why three copies} \code{potrf} overwrites its matrix, \code{trsm}
and \code{trsv} their right-hand sides. \(M\), \(X\) and \(y\) are
\code{ulm_lam.in}, so \code{ulm-lam-results-to-out-params} gave them
\code{bufferization.writable = false}
(\srcref{lib/UlmLamResultsToOutParams.cpp}{344}{368}), and One-Shot
Bufferize must not write into them: it copies \(M\) and \(y\). For \(X\)
the copy was already in the IR (\code{gecopy}, a transposing copy that
\code{trsm} then overwrites), so no further copy is needed. With
\code{InOut(...)} parameters the caller allows overwriting;
\code{examples/gls/check.sh} verifies that \code{potrf} then works on the
caller's \(M\) and only \(X\) is copied.

\paragraph{One value, two uses, one buffer} \(W\) lives in \code{\%view_2},
written once by \code{trsm} and read by \code{syrk} and \code{gemv}. SSA
use counts are not copies: a value with two uses needs a second buffer only
if one use would overwrite it while the other still reads it.

\paragraph{The result} \(\beta\) is computed directly in the caller's
output buffer \code{\%arg5}: \code{gemv} writes it, the two \code{trsv}
update it in place.

\paragraph{The workspace} All temporaries are views of one byte buffer
\code{\%25} at offsets computed from \(n\) and \(m\), each rounded up to 64
bytes (\cref{fig:workspace}). The caller can pass that buffer, query its size
first (LAPACK's \code{lwork} convention), or let the function allocate it.

\paragraph{\code{info}} Each factorization is a numbered check. If
\code{potrf} reports failure, the rest of the block does not run and the
function returns the check's number. In C++ this is the return value of
\code{gls}; in Julia, an \code{InfoError}.

\begin{theory}[ (\FGs{10.3}, \FGs{2.7}, \FGs{1.4})]
Tensors have value semantics, memrefs reference semantics,
\FG{2.7}{Value semantics vs reference/memory semantics in IRs}.
Bufferization chooses, for every DPS result, between reusing the
destination's buffer and copying, \FG{10.3}{Bufferization and
destination-passing style}. The in-place condition (no other live reader of
the old value, and the buffer may be written) is a uniqueness condition, as
\FG{1.4}{Uniqueness types vs linear types} explains; here it is decided by
an analysis, not by a type system. \Cref{ch:storage} goes into detail.
\end{theory}

% ---------------------------------------------------------------------------
\section{The kernel sequence in C++}
\label{sec:gls-kernels}

\code{convert-ulm-lamk-to-ulmblas} lowers the memref program to the EmitC
dialect, and \code{mlir-translate --mlir-to-cpp} prints C++. With the
Julia backend the calls go to CBLAS (the C interface to BLAS) and LAPACK
through Julia's libblastrampoline, which forwards them to OpenBLAS here. The
calls of \texttt{7-gls.cpp}, in order (\code{v1} \(=n\), \code{v2} \(=m\),
\code{v3} \(=X\T\) row-major, \code{v4} \(=M\), \code{v5} \(=y\),
\code{v6} \(=\beta\)):

\begin{lstlisting}[style=cpp]
ulm_lam_copy(v31, v4, v32);                                          // M -> workspace
int64_t v33 = ulm_lam_cblas::potrf<std::ptrdiff_t>(true, v1, v31, v1, 1);
ulm_lam_gecopy(v1, v2, v3, 1, v1, v41, v2, 1);                       // X (transposing copy)
(void)ulm_lam_cblas::trtrs<std::ptrdiff_t>(
    true, false, false, false, v1, v2, v31, v1, 1, v41, v2, 1);      // W = L1^-1 X
ulm_lam_cblas::sylrk<std::ptrdiff_t>(v2, v1, v9, v41, 1, v2,
                                     v10, v39, v2, 1);               // S = W^T W (lower)
int64_t v42 = ulm_lam_cblas::potrf<std::ptrdiff_t>(true, v2, v39, v2, 1);
ulm_lam_copy(v45, v5, v46);                                          // y -> workspace
ulm_lam_cblas::trlsv<std::ptrdiff_t>(v1, false, v31, v1, 1, v45, 1);  // z = L1^-1 y
ulm_lam_cblas::gemv<std::ptrdiff_t>(v2, v1, v9, v41, 1, v2,
                                    v45, 1, v10, v6, 1);             // beta = W^T z
ulm_lam_cblas::trlsv<std::ptrdiff_t>(v2, false, v39, v2, 1, v6, 1);   // beta = L2^-1 beta
ulm_lam_cblas::trusv<std::ptrdiff_t>(v2, false, v39, 1, v2, v6, 1);   // beta = L2^-T beta
\end{lstlisting}

Two details show what lowering does with the remaining structure. The
transposition flags are gone: a transposed operand is the same pointer with
swapped strides (\code{v39, 1, v2} versus \code{v39, v2, 1}), so
\(L_2\T\) is passed as an \emph{upper} triangular matrix to
\code{trusv}. And the matrix properties are gone entirely; what is left of
them is the choice of routines and their \code{uplo} arguments.

\begin{center}\small
\begin{tabularx}{\linewidth}{@{}lllL@{}}
\toprule
\textbf{Kernel} & \textbf{Library} & \textbf{Flops} & \textbf{Computes} \\
\midrule
\code{potrf} & LAPACK & \(n^3/3\) & \(M = L_1 L_1\T\), in a copy of \(M\) \\
\code{trsm} (\code{trtrs}) & BLAS 3 & \(n^2 m\) & \(W = L_1\inv X\), in a copy of \(X\) \\
\code{syrk} & BLAS 3 & \(n m^2\) & \(S = W\T W\), lower triangle \\
\code{potrf} & LAPACK & \(m^3/3\) & \(S = L_2 L_2\T\), in place \\
\code{trsv} & BLAS 2 & \(n^2\) & \(z = L_1\inv y\), in a copy of \(y\) \\
\code{gemv} & BLAS 2 & \(2nm\) & \(\beta \leftarrow W\T z\) \\
\code{trsv}, \code{trsv} & BLAS 2 & \(2m^2\) & \(\beta \leftarrow L_2\invT L_2\inv \beta\) \\
\bottomrule
\end{tabularx}
\end{center}

% ---------------------------------------------------------------------------
\section{A numerical example}
\label{sec:gls-numbers}

Small inputs make every step checkable by hand. Take \(n=4\), \(m=2\),
\[
X=\begin{pmatrix}1&0\\1&1\\1&2\\1&3\end{pmatrix},\quad
M=\begin{pmatrix}4&2&0&0\\2&5&2&0\\0&2&5&2\\0&0&2&5\end{pmatrix},\quad
y=\begin{pmatrix}1\\2\\4\\4\end{pmatrix}.
\]
\(M\) is SPD (eigenvalues between 1.58 and 8.13) with Cholesky factor
\(L_1\) = 2 on the diagonal and 1 below it, so every step up to \(S\) is
exact in binary floating point:
\[
W = L_1\inv X = \begin{pmatrix}1/2&0\\1/4&1/2\\3/8&3/4\\5/16&9/8\end{pmatrix},\quad
S = W\T W = \begin{pmatrix}141/256 & 97/128\\ 97/128 & 133/64\end{pmatrix},\quad
X\T M\inv y = \begin{pmatrix}363/256\\375/128\end{pmatrix}.
\]

\paragraph{Expected} In exact rational arithmetic,
\[
  \beta = \Bigl(\frac{93}{73},\ \frac{69}{73}\Bigr) = (1.273972602739726\ldots,\ 0.945205479452054\ldots).
\]

\paragraph{Obtained} \code{examples/gls.jl} prints:

\begin{textout}
==== exact β (rational): Rational{BigInt}[93//73, 69//73]
==== exact β (Float64):  [1.273972602739726, 0.9452054794520548]
==== ulm.lam β: [1.2739726027397262, 0.9452054794520546]
==== Julia expert  (X'*(M\X)) \ (X'*(M\y)): [1.2739726027397262, 0.9452054794520547]
==== Julia formula inv(X'*inv(M)*X)*X'*inv(M)*y: [1.2739726027397262, 0.9452054794520547]
==== max abs error vs exact: 2.220446049250313e-16
\end{textout}

ulm.lam's result differs from the correctly rounded \(\beta\) by one unit
in the last place in each component, and from Julia's two evaluations in the
second component by one unit. All three are correct to working accuracy;
they differ because they round differently, not because one is wrong.

\paragraph{Failing and rejected inputs} The same program with inputs that
violate its assumptions (all from \code{outputs/gls.out}):

\begin{center}\small
\begin{tabularx}{\linewidth}{@{}P{4.6cm}P{5.3cm}L@{}}
\toprule
\textbf{Input} & \textbf{Result} & \textbf{Origin of the check} \\
\midrule
\(M_{33} = -1\): symmetric, indefinite & \code{InfoError: gls: info 1: Cholesky factorisation of v4: not positive definite} & \Orun: \code{potrf} of \(M\) \\
\(X\) with two equal columns, declared \code{full_rank} & \code{InfoError: gls: info 2: Cholesky factorisation of a computed matrix: not positive definite} & \Orun: \code{potrf} of \(S\) \\
\(y\) of length 3 & \code{LamError: gls: y: size :n = 3, but 4 from an earlier argument} & \Orun: the Julia runtime compares sizes before the call \\
\(M_{12} = 100\) (Julia's upper triangle), declared \code{spd} & \code{InfoError: ... info 1 ...} & \Orun, by luck: \code{potrf} reads the triangle that was changed \\
\(M_{21} = 100\) (Julia's lower triangle), declared \code{spd} & \(\beta\) of \code{Symmetric(M, :U)}; no error & \Onone: symmetry is never checked, one triangle is ignored \\
\code{Mat(:n,:m) * Mat(:n,:m)} in a \code{@lam} body & \code{LamError: product of nxm and nxm: m columns but n rows} & at trace time, by the frontend (\cref{ch:shapes}) \\
\bottomrule
\end{tabularx}
\end{center}

Two remarks. The messages name \code{v4}, the fourth C parameter, where
\code{@code_lam} says \code{M}: the description comes from the C++
interface comment. And the rank-deficient \(X\) is caught only because
\(S\) became exactly singular here. A nearly rank-deficient \(X\) can pass
the factorization and give an inaccurate \(\beta\), or fail although \(X\)
has full rank (\cref{sec:kernels-fp}).

% ---------------------------------------------------------------------------
\section{The chapter in one table}
\label{sec:gls-summary}

\begin{center}\small
\begin{tabularx}{\linewidth}{@{}P{3.1cm}P{4.6cm}P{3.3cm}L@{}}
\toprule
\textbf{Step} & \textbf{MLIR mechanism} & \textbf{Theory} & \textbf{FG} \\
\midrule
tracing & text emitted by overloaded Julia operators & tracing eDSL & \FGs{3.7} \\
facts on matrices & builtin tensor type with an attribute in its encoding & refinement types & \FGs{1.7} \\
normal form & \code{OpRewritePattern}s, greedy driver & rewriting, termination & \FGs{3.4}, \FGs{3.5} \\
\(X\T M\inv X\) is SPD & \code{InferTypeOpInterface} per op & type inference; abstract interpretation & \FGs{1.10}, \FGs{5.2} \\
kernel selection & custom pass, DP over leaves & instruction selection & \FGs{4.5}, \FGs{3.3} \\
copies, in place & DPS ops, One-Shot Bufferize & value vs reference semantics; uniqueness & \FGs{10.3}, \FGs{2.7}, \FGs{1.4} \\
failed factorization & \code{info} results, \code{scf.if} & run-time check & --- \\
C++ & dialect conversion to EmitC & lowering, legalization & \FGs{3.3}, \FGs{4.6} \\
\bottomrule
\end{tabularx}
\end{center}
